% -------------------------------------------------------------------------
% WBS ID: RES-MOD-IBC
% Deliverable: Initial, Boundary and Interface Condition Models
% Status: Draft complete
% Evidence status: Research specification complete for regional NLSWE and local 3D boundary roles
% Review status: Not reviewed
% Last updated: 2026-07-18
% Dependencies: RES-MOD-EQS, RES-DAT-GEO, RES-DAT-SRC, RES-MOD-MUL
% Downstream interfaces: RES-NUM-IBC, RES-NUM-WET, RES-VER-SPEC
% -------------------------------------------------------------------------

\section{RES-MOD-IBC --- Initial, Boundary and Interface Condition Models}
\label{sec:res-mod-ibc}

\subsection{Purpose and Scope}
\label{sec:res-mod-ibc-purpose}

% TODO: Define what this Deliverable covers and explicitly excludes.

\subsection{Research Questions}
\label{sec:res-mod-ibc-research-questions}

% TODO: State the questions that must be answered before the Deliverable
% can be accepted.

\subsection{Terminology and Definitions}
\label{sec:res-mod-ibc-definitions}

% TODO: Define the technical terminology, symbols, quantities and
% classifications used in this Deliverable.

\subsection{Literature Evidence}
\label{sec:res-mod-ibc-literature-evidence}

% TODO: Record evidence from peer-reviewed papers, authoritative datasets,
% standards, technical reports and primary documentation.
%
% Do not create a detached literature-summary list. Explain how each source
% supports, contradicts or limits a project decision.

\textcite{higuera2013realistic} is retained as evidence that
Navier--Stokes numerical wave tanks require compatible wave generation,
active absorption and relaxation/open-boundary treatments. Source
role: `3D-WAVEBC`. Project conclusion: the local inlet, outlet and
lateral ocean cuts shall be treated as wave-boundary components rather
than generic zero-gradient faces. Limitation: exact relaxation-zone
lengths and damping profiles remain numerical and validation choices.

\textcite{takase2016hybrid} and \textcite{romano2025coupling} support
independent regional and local domains connected through interface
state transfer. Source role: `3D-COUPLING`. Project conclusion: local
inlet histories shall be reconstructable from regional water depth,
surface elevation and discharge histories. Limitation: the initial
depth-uniform hydrostatic lift is not a unique reconstruction of the
three-dimensional state.

\subsection{Governing Relations and Quantitative Definitions}
\label{sec:res-mod-ibc-governing-relations}

% TODO: Record relevant equations, dimensionless groups, empirical models,
% extraction rules or quantitative definitions.
%
% Use align environments for displayed equations.
% Every displayed equation must have a unique \label{eq:...}.
% Do not place punctuation after displayed equations.

\subsection{System Boundary}
\label{sec:res-mod-ibc-system-boundary}

% TODO: Complete this RES-MOD Domain-specific subsection.

\subsection{State Variables and Parameters}
\label{sec:res-mod-ibc-state-variables-parameters}

% TODO: Complete this RES-MOD Domain-specific subsection.

\subsection{Continuous Governing Model}
\label{sec:res-mod-ibc-continuous-governing-model}

% TODO: Complete this RES-MOD Domain-specific subsection.

\subsection{Source Terms and Closures}
\label{sec:res-mod-ibc-source-terms-closures}

% TODO: Complete this RES-MOD Domain-specific subsection.

\subsection{Initial, Boundary and Interface Conditions}
\label{sec:res-mod-ibc-initial-boundary-interface-conditions}

The regional two-dimensional model uses a horizontal depth-averaged
state. Initial conditions therefore consist of the cellwise water depth,
depth-integrated discharges, static bed, and any time-dependent
earthquake bed increments. Cross-sectional views are diagnostic outputs
derived from the horizontal solution and shall not be used as separate
one-dimensional solved domains.

Artificial ocean boundaries are represented by boundary states
$\mathbf{U}_{\mathrm{B}}$ supplied to the numerical flux. The baseline
mathematical boundary roles are:

\begin{table}[htbp]
  \centering
  \small
  \renewcommand{\arraystretch}{1.2}
  \begin{tabular}{
      p{0.26\textwidth}
      p{0.30\textwidth}
      p{0.34\textwidth}
    }
    \hline
    Boundary or interface
    &
    Regional treatment
    &
    Status
    \\
    \hline
    Offshore artificial boundary
    &
    Characteristic or radiation open boundary
    &
    Established at specification level; validation pending
    \\
    Lateral artificial ocean cuts
    &
    Radiation or open boundary, with sponge damping where required
    &
    Established at specification level; domain-specific tuning pending
    \\
    Physical coastline
    &
    Topography plus wetting--drying, not an external wall boundary
    &
    Established at specification level
    \\
    True wall or resolved structure
    &
    Reflective boundary state
    &
    Established for walls; final defence representation remains
    RES-GEO/RES-MOD-MUL work
    \\
    Regional-to-local interface
    &
    Conservative transfer of $h$, $\eta$, $q_x$, $q_y$ and derived
    boundary histories
    &
    Coupling algorithm deferred
    \\
    \hline
  \end{tabular}
  \caption{Regional boundary and interface condition roles}
  \label{tab:res-mod-ibc-regional-boundary-roles}
\end{table}

\subsubsection{Local Three-Dimensional Boundary Roles}
\label{sec:res-mod-ibc-local-boundary-roles}

The local boundary partition is defined in
\cref{eq:res-mod-sys-local-boundary-partition}. The continuous local
boundary roles are:

\begin{table}[htbp]
  \centering
  \small
  \renewcommand{\arraystretch}{1.2}
  \begin{tabular}{
      p{0.20\textwidth}
      p{0.31\textwidth}
      p{0.38\textwidth}
    }
    \hline
    Boundary
    &
    Continuous role
    &
    Project status
    \\
    \hline
    $\Gamma_{\mathrm{in}}$
    &
    Regional-to-local incident-wave forcing
    &
    Selected for one-way replay; also used by simultaneous coupling
    \\
    $\Gamma_{\mathrm{out}}$
    &
    Downstream open or radiation boundary with absorption where needed
    &
    Selected as a boundary role; exact formula and relaxation length
    remain numerical choices
    \\
    $\Gamma_{\mathrm{side}}$
    &
    Lateral open-ocean cuts
    &
    Selected as open/radiation boundaries, not physical walls
    \\
    $\Gamma_{\mathrm{top}}$
    &
    Open atmosphere
    &
    Selected as an open pressure/reference-air boundary
    \\
    $\Gamma_{\mathrm{T}}$
    &
    Fixed terrain and bathymetry
    &
    Selected as a no-slip wall with wall-model details owned by
    `RES-NUM-IBC`
    \\
    $\Gamma_{\mathrm{B}}$
    &
    Fixed rigid barrier surface
    &
    Selected for baseline hydrodynamics; structural motion deferred
    to activated FSI
    \\
    \hline
  \end{tabular}
  \caption{Continuous local three-dimensional boundary-condition roles}
  \label{tab:res-mod-ibc-local-boundary-roles}
\end{table}

For one-way replay the regional model supplies
$h_{\mathrm{R}}$, $\eta_{\mathrm{R}}$, $q_{x,\mathrm{R}}$ and
$q_{y,\mathrm{R}}$ on the interface. The depth-averaged velocities are:

\begin{align}
  u_{\mathrm{R}}
  &=
  \frac{
    q_{x,\mathrm{R}}
  }{
    h_{\mathrm{R}}
  }
  \label{eq:res-mod-ibc-regional-to-local-u}
  \\
  v_{\mathrm{R}}
  &=
  \frac{
    q_{y,\mathrm{R}}
  }{
    h_{\mathrm{R}}
  }
  \label{eq:res-mod-ibc-regional-to-local-v}
\end{align}

The baseline local inlet reconstruction is:

\begin{align}
  \alpha_{\mathrm{in}}
  &=
  \mathcal{H}
  \left(
    \eta_{\mathrm{R}}-z
  \right)
  \label{eq:res-mod-ibc-local-inlet-alpha}
  \\
  \mathbf{u}_{\mathrm{in}}
  &=
  \left[
    u_{\mathrm{R}},\,
    v_{\mathrm{R}},\,
    0
  \right]^{\mathsf{T}}
  \label{eq:res-mod-ibc-local-inlet-velocity}
  \\
  p_{\mathrm{in}}
  &=
  \rho_{\mathrm{w}} g
  \max
  \left(
    \eta_{\mathrm{R}}-z,\,
    0
  \right)
  \label{eq:res-mod-ibc-local-inlet-pressure}
\end{align}

where $\mathcal{H}$ is a phase-indicator reconstruction. This is a
depth-uniform, hydrostatic lift and is valid only where the interface
is upstream of strong local vertical acceleration, breaking and
barrier-induced separation.

On the fixed terrain and barrier surfaces the baseline velocity
condition is:

\begin{align}
  \mathbf{u}_{\mathrm{f}}
  &=
  \mathbf{0}
  \qquad
  \mathrm{on}
  \quad
  \Gamma_{\mathrm{T}}
  \cup
  \Gamma_{\mathrm{B}}
  \label{eq:res-mod-ibc-local-no-slip}
\end{align}

The no-slip condition defines the continuous wall role. The high-
Reynolds-number wall-function approximation used by the discrete model
belongs to `RES-NUM-IBC` and shall not be mistaken for a change in the
continuous boundary condition.

\subsection{Model Coupling Requirements}
\label{sec:res-mod-ibc-model-coupling-requirements}

% TODO: Complete this RES-MOD Domain-specific subsection.

The local boundary histories shall remain event-independent in
formulation. A Tohoku case, or any later event, changes the supplied
time series and terrain data, not the mathematical boundary roles in
\cref{tab:res-mod-ibc-local-boundary-roles}.

\subsection{Sufficient Conditions for Validity}
\label{sec:res-mod-ibc-sufficient-conditions-validity}

% TODO: Complete this RES-MOD Domain-specific subsection.

\subsection{Candidate Approaches}
\label{sec:res-mod-ibc-candidate-approaches}

% TODO: Define the credible candidate approaches considered.

\subsection{Critical Comparison}
\label{sec:res-mod-ibc-critical-comparison}

% TODO: Compare candidates using physically and project-relevant criteria.
% Do not select an approach solely on popularity or implementation ease.

\subsection{Assumptions and Applicability}
\label{sec:res-mod-ibc-assumptions}

% TODO: State assumptions and the conditions under which they are valid.

\subsection{Limitations and Failure Conditions}
\label{sec:res-mod-ibc-limitations}

% TODO: State where the evidence, model or selected approach ceases to be
% reliable or applicable.

\subsection{Project-Specific Conclusions}
\label{sec:res-mod-ibc-conclusions}

% TODO: Convert the evidence into conclusions specific to the Studentship
% project. Do not merely restate literature findings.

\subsection{Selected Approach and Rationale}
\label{sec:res-mod-ibc-selected-approach}

% TODO: Record the selected approach only after sufficient evidence exists.
% State the alternatives rejected and the reasons for rejection.

The selected regional condition model is a face-state boundary
formulation. Boundary conditions act during every residual evaluation,
not as a post-step correction. This preserves consistency with the
finite-volume conservation update and lets open, reflective, sponge and
interface boundaries share one mathematical mechanism.

\subsection{Unresolved Decisions}
\label{sec:res-mod-ibc-unresolved-decisions}

% TODO: Record unresolved questions, required evidence and decision owners.

The regional boundary-condition roles are established, but the following
remain open:

\begin{itemize}
  \item final QGIS/geospatial domain lock and artificial-boundary
    placement;
  \item sponge-layer width and damping scale after domain selection;
  \item final regional-to-local transfer operators;
  \item final representation of documented defences as topography,
    wall faces, holes, porous interfaces, or local 3D extraction
    regions.
  \item local inlet relaxation treatment and pressure compatibility
    under strong reflection;
  \item local wall-function and roughness parameters for terrain and
    barrier surfaces;
  \item local outlet and side-boundary absorption settings.
\end{itemize}

\subsection{Requirements and Handoffs}
\label{sec:res-mod-ibc-requirements}

% TODO: State requirements passed to other Research Domains, Software,
% verification activities or later execution work.

`RES-NUM-IBC` shall express each boundary role as a discrete boundary
state or transfer operator. `RES-DAT-GEO` shall provide the final
domain, coastline, datum and coordinate transforms before boundary
locations are locked.

`RES-VER-FRM` and `RES-VER-MET` shall test whether boundary reflections
or inlet reconstruction errors alter run-up, overtopping, pressure,
force, moment or candidate ranking.

\subsection{Deliverable Completion Record}
\label{sec:res-mod-ibc-completion-record}

% TODO: Record whether the Deliverable is:
%
% - Not started
% - In progress
% - Draft complete
% - Reviewed
% - Accepted
%
% Acceptance requires:
%
% - the research questions to be answered;
% - claims to be supported by appropriate evidence;
% - assumptions and limitations to be explicit;
% - the project-specific conclusion to be stated;
% - downstream requirements to be traceable;
% - unresolved decisions to be closed or formally deferred.
